% \section{Problem reformulation}
% \label{Section: Problem Reformulation}
% \subsection{W-type Regularized Deep Hashing Model}
% Deep hashing methods aim to learn network parameters \(\bx \in \calX \subseteq \RR^{d_{\bx}}\) to generate discrete codes that preserve semantic similarity effectively in the Hamming space. 
% A natural approach is to minimize a similarity-preserving loss function \( \frac{1}{n}\sum_{i=1}^n L_s(\text{sgn}(\bC_i(\bx)))\) where $n$ is the number of training samples, $\bC_i(\bx)$ represents the \(i\)-th network output and \(f_i\) is the similarity-preserving loss function (e.g., pairwise cross-entropy loss \cite{zhuDeepHashingNetwork2016}). However, the function \( sgn(\cdot) \) renders standard backpropagation infeasible due to its zero gradient for all nonzero inputs. This discreteness poses a central challenge in deep hashing.

% A W-type regularization (e.g., $\lambda | |z| - 1 |$) has effectively mitigated this issue. This type of regularization encourages outputs to approach binary values \(\{-1, 1\}\) during training. Specifically, the model is trained to produce continuous outputs \(\bu_i = \bD_i(\bx)\) with W-type regularization, and binary codes are obtained post-training via \(\bb_i = \text{sgn}(\bu_i)\).  Consequently, we formulate the core optimization problem for W-type regularized deep hashing methods as follows,
% \begin{align}
% \label{Eq: the primal deep hashing problem}
% \min_{\bx \in \calX} \quad \frac{1}{n} \sumn f_i (\bx) + \sumn h(\bD_i(\bx)), 
% \end{align}where \(\frac{1}{n}\sumn f_i\) is the similarity-preserving loss, \(D_i\) is the \(i\)-th network output and \(h\) is a W-type regularization.
% % \footnote{In the problem \eqref{Eq: the primal deep hashing problem}, we use vectorized \(\bz\), thus the regularization is also \(\| |\bz| - \bm{1}\|_1 = \sumn | |z_i| - 1 |\)}.
% \subsection{A Two-Block Structured Reformulation}
% We then employ variable splitting % \cite{wangLiftedBregmanTraining2022, zengBlockCoordinateDescent2018, lauProximalBlockCoordinate2018}, ]\
% to reformulate problem \eqref{Eq: the primal deep hashing problem} equivalently as the constrained problem:
% \begin{align} 
% \label{Eq: constrained_dlh} 
% \min_{\bx \in \mathcal{X}, \bB \in \mathcal{B}}&\quad \frac{1}{n} \sumn f_i (\bx) + \sum_{i=1}^n h(\bb_i) \notag \\ 
% \text{s.t.} &\quad \bb_i = \bD_i(\bx), \; \forall i=1,\dots,n. 
% \end{align} Combined with the quadratic penalty method \cite{nocedal1999numerical}, this transformation yields the following two-block finite-sum optimization problem:
% \begin{align} 
% \label{Prop: block-wise regularized problem} 
% \nonumber&\min_{\bx \in \mathcal{X}, \bB \in \mathcal{B}} J(\bx, \bB) \\
% &= \overbrace{\underbrace{\frac{1}{n}\sum_{i=1}^n f_i(\bx)}_{f(\bx)} + \underbrace{\frac{\gamma}{2n} \sum_{i=1}^n \| \bD_i(\bx) - \bb_i\|^2}_{P(\bx, \bB)}}^{F(\bx, \bB) = \frac{1}{n}\sumn F_i(\bx, \bb_i)} + \underbrace{\sum_{i=1}^n h(\bb_i)}_{h(\bB)}. \tag{P}
% \end{align} Here, \( \bB = [\bb_1, \bb_2, \dots, \bb_n]\) and \( \gamma > 0 \) is the penalty parameter. This reformulation decomposes the original problem \eqref{Eq: the primal deep hashing problem} into simpler structural blocks, enabling us to exploit the problem structure better. Building upon this problem, we develop a stochastic primal-dual algorithm that effectively addresses the challenges above.
\section{Problem reformulation}
\label{Section: Problem Reformulation}
\subsection{W-type Regularized Deep Hashing Model}
Deep hashing methods aim to learn network parameters \(\bx \in \calX \subseteq \RR^{d_{\bx}}\) to generate discrete codes that preserve semantic similarity effectively in the Hamming space. 
A natural approach is to minimize a similarity-preserving loss function \( \frac{1}{n}\sum_{i=1}^n L_s(\text{sgn}(\bC_i(\bx)))\) where $n$ is the number of training samples, $\bC_i(\bx)$ represents the \(i\)-th network output and \(f_i\) is the similarity-preserving loss function (e.g., pairwise cross-entropy loss \cite{zhuDeepHashingNetwork2016}). However, the function \( sgn(\cdot) \) renders standard backpropagation infeasible due to its zero gradient for all nonzero inputs. This discreteness poses a central challenge in deep hashing.

A W-type regularization (e.g., $\lambda | |z| - 1 |$) has effectively mitigated this issue. This type of regularization encourages outputs to approach binary values \(\{-1, 1\}\) during training. Specifically, the model is trained to produce continuous outputs \(\bu_i = \bD_i(\bx)\) with W-type regularization, and binary codes are obtained post-training via \(\bb_i = \text{sgn}(\bu_i)\).  Consequently, we formulate the core optimization problem for W-type regularized deep hashing methods as follows,
\begin{align}
    \label{Eq: the primal deep hashing problem}
    \min_{\bx \in \calX} \quad \frac{1}{n} \sumn f_i (\bx) + \sumn h(\bD_i(\bx)),
\end{align}
where \(\frac{1}{n}\sumn f_i\) is the similarity-preserving loss, \(D_i\) is the \(i\)-th network output and \(h\) is a W-type regularization.
\subsection{A Two-Block Structured Reformulation}
We then employ variable splitting to reformulate problem \eqref{Eq: the primal deep hashing problem} equivalently as the constrained problem:
\begin{align}
    % \label{Eq: constrained_dlh}
    \min_{\bx \in \mathcal{X}, \bB \in \mathcal{B}} \quad & \frac{1}{n} \sumn f_i (\bx) + \sum_{i=1}^n h(\bb_i) \notag \\
    \text{s.t.} \quad & \bb_i = \bD_i(\bx), \; \forall i=1,\dots,n. \notag
\end{align}
Combined with the quadratic penalty method \cite{nocedal1999numerical}, this transformation yields the following two-block finite-sum optimization problem:
\begin{align}
    \label{Prop: block-wise regularized problem}
    \min_{\bx \in \mathcal{X}, \bB \in \mathcal{B}} J(\bx, \bB) = \overbrace{\underbrace{\frac{1}{n}\sum_{i=1}^n f_i(\bx)}_{f(\bx)} + \underbrace{\frac{\gamma}{2n} \sum_{i=1}^n \| \bD_i(\bx) - \bb_i\|^2}_{P(\bx, \bB)}}^{F(\bx, \bB) = \frac{1}{n}\sumn F_i(\bx, \bb_i)} + \underbrace{\sum_{i=1}^n h(\bb_i)}_{h(\bB)}. \tag{P}
\end{align}
Here, \( \bB = [\bb_1, \bb_2, \dots, \bb_n]\) and \( \gamma > 0 \) is the penalty parameter. This reformulation decomposes the original problem \eqref{Eq: the primal deep hashing problem} into simpler structural blocks, enabling us to exploit the problem structure better. 
Building upon Problem~\eqref{Prop: block-wise regularized problem}, we develop a stochastic primal-dual algorithm \blue{that addresses the lack of tractable proximal mappings in the nonsmooth composite form and the absence of convergence guarantees in existing deep hashing methods.}
